%% kapa: exam=1 round=21 year=2010 period=2 form=A subject=appraisal_law subject_ko=부동산관계법규 seq=10 number=50
%% source: Q-Net 감정평가사 기출문제 (cst00309) artlSeq=5024181 file=extracted/1008/5024181/2010년 제21회 감정평가사 1차시험 기출문제/2교시.hwp
%% conversion: kordoc 4.24.51 -> tools/gen1.py (auto); review=proofread-gpt-sol; flags=none
\begin{question}{50}
국토의 계획 및 이용에 관한 법령상 토지거래계약에 관한 허가구역에 관한 설명으로 옳지 않은 것은?
\begin{choices}
  \item 허가구역의 지정권자는 국토해양부장관이다.
  \item 허가구역의 지정은 허가구역의 지정을 공고한 날로부터 5일 후에 그 효력이 발생한다.
  \item 허가구역은 5년 이내의 기간을 정하여 지정할 수 있다.
  \item 지가가 급격히 상승하는 지역으로서 법령에 의한 개발사업이 진행중이거나 예정되어 있는 지역은 허가구역으로 지정할 수 있다.
  \item 「민사집행법」에 따른 경매의 경우에도 경매대상물인 토지가 허가구역 안에 있는 경우에는 토지거래계약허가를 받아야 한다.
\end{choices}
\end{question}
\answer{5}
